\documentclass[10pt]{article}

% Packages
\usepackage[utf8]{inputenc}
\usepackage[T1]{fontenc}
\usepackage{amsmath}
\usepackage{graphicx}
\usepackage[margin=1in]{geometry}
\usepackage[numbers,sort&compress]{natbib}
\usepackage{hyperref}
\usepackage{lineno}
\usepackage{setspace}
\usepackage{url}
\usepackage{booktabs}

% Nature Medicine: article titles omitted from references
\bibliographystyle{naturemag}

% Line numbers for review
\linenumbers
\doublespacing

% ============================================================
\title{The threat of synthetic respondents extends to clinical mental health screening}

\author{%
  Kiant\'{e} Fernandez$^{1}$ \and Laura A. Berner$^{2}$ \and Blair R. K. Shevlin$^{2}$
}

\date{}

\begin{document}

\maketitle

\begin{center}
  $^{1}$Department of Psychology, University of California, Los Angeles, Los Angeles, California \\
  $^{2}$Center for Computational Psychiatry, Department of Psychiatry, Icahn School of Medicine at Mount Sinai, New York, New York \\
\end{center}

\vspace{1em}

% ============================================================
% Abstract (up to 150 words, unreferenced)
% ============================================================
\begin{abstract}
Online platforms have enabled clinical mental health research at unprecedented scale, yet large language models (LLMs) now threaten the integrity of remotely collected data. Recent work has demonstrated that LLM-based synthetic respondents can complete general surveys with human-like plausibility, but the vulnerability of validated clinical psychiatric instruments remains unexamined. Here, we show that a commercially available LLM, given only brief diagnostic persona descriptions, produces clinically differentiated responses across seven standardized psychiatric screening instruments spanning mood, anxiety, obsessive-compulsive, trauma, psychotic, neurodegenerative, and eating disorders. In a simulation of 2,106 synthetic personas, diagnosis-congruent personas generated scores exceeding clinical cutoffs on five of seven instruments, with scores scaling monotonically with assigned severity. This mimicry required no specialized configuration. Coherent symptom endorsement can no longer serve as a proxy for authentic participation, urgently necessitating standardized verification methods for clinical online data.
\end{abstract}

\newpage

% ============================================================
% Main Text (up to 2,000 words total incl. abstract, refs, legends)
% NO HEADINGS per Nature Medicine Brief Communication format
% ============================================================

Online platforms and crowdsourcing services have transformed how clinical researchers study mental health\cite{gillan2021}. Surveys administered through services such as Amazon Mechanical Turk and Prolific enable rapid, low-cost access to large and diverse samples, including populations that are difficult to reach through traditional recruitment due to stigma, geographic barriers, or low base rates\cite{chandler2016}. This infrastructure has become particularly important for research on psychiatric conditions, where validated screening instruments commonly serve as gatekeepers for study eligibility and primary outcome measures. However, the assumption that coherent, contextually appropriate survey responses require a human respondent is increasingly questioned.

Westwood recently demonstrated that autonomous LLM-based respondents can complete online surveys with high fidelity\cite{westwood2025}. These LLM-based agents passed attention checks, produced consistent psychometric profiles, and evaded standard detection methods. Critically, each completed survey cost approximately \$0.005, making large-scale fabrication economically trivial. Moreover, the average user can now employ semi-autonomous AI tools with agents built into browsers that can take surveys with little input\cite{ning_survey_2025}. While the current scale of AI contamination on online research platforms remains uncertain, these initial probes demonstrate that the threat to online research is broad and deepening\cite{panizza_how_2026}. Yet to our knowledge, none of these studies have examined whether LLMs can produce plausible responses on the clinical instruments that mental health researchers regularly rely on to screen, characterize, and study psychiatric populations.

Clinical psychiatric instruments may be especially vulnerable to LLM mimicry for several reasons. First, mental health content, including descriptions of symptom presentations, diagnostic criteria, and self-report questionnaire items, is extensively represented in the text corpora used to train these models\cite{Jin2025ApplicationsLLMsMentalHealth}. For example, word embeddings trained exclusively on LLM-generated synthetic self-reports can predict mental health screening scores with high accuracy\cite{moell2026}, suggesting that even synthetic clinical text may capture patterns associated with diagnostic information. Second, many clinical studies recruit participants online precisely because the populations of interest are difficult to access through other means\cite{chandler2016}. Third, researchers routinely rely on published cutoff scores to validate study eligibility, treating coherent symptom endorsement as implicit proof of a genuine respondent. An agent need only exceed a known threshold to infiltrate a clinical sample, turning quality-control metrics into a vulnerability.

To evaluate this vulnerability, we conducted a simulation study in which a commercially available LLM (Google Gemini 2.0 Flash) was used to generate responses to seven validated clinical instruments. We created 2,106 synthetic personas from 13 DSM-informed clinical diagnostic profiles, three severity levels (mild, moderate, severe), and demographic variables. Diagnostic profiles spanned mood disorders, anxiety disorders, obsessive-compulsive disorders, psychotic disorders, eating disorders, and neurodegenerative conditions. The model received a straightforward system prompt describing the persona's clinical profile and demographic background. No adversarial prompt engineering or fine-tuning was employed, and the prompt contained no information about instrument content, scoring rubrics, or clinical cutoffs.

LLM-generated personas produced clinically differentiated response patterns across all seven instruments (Fig.~\ref{fig:1}; Table~\ref{tab:main-effects}). We note two key findings. First, diagnostic specificity was instrument-dependent. On disorder-specific instruments, target personas scored substantially higher than nontarget controls: OCI-R ($b = 15.76$, 95\% CI [12.20, 19.31], $P < .001$), PCL-5 ($b = 14.63$, 95\% CI [10.09, 19.17], $P < .001$), MDQ ($b = 9.24$, 95\% CI [8.42, 10.07], $P < .001$), PQ-16 ($b = 6.15$, 95\% CI [3.02, 9.28], $P < .001$), and EDE-Q ($b = 2.95$, 95\% CI [2.26, 3.64], $P < .001$). By contrast, the PHQ-9 and GAD-7 showed no diagnostic specificity: target and nontarget personas produced nearly identical mean scores (PHQ-9: 13.58 vs 13.45, $P = .915$; GAD-7: 10.78 vs 11.54, $P = .306$). Although the model generated elevated depression and anxiety symptoms across all diagnostic profiles, this is consistent with their known comorbidities with other disorders\cite{McGrath._2020}. Second, despite this specificity gap, scores increased monotonically with assigned severity within target groups on all seven instruments (all $P < .001$; Table~\ref{tab:main-effects}). On the PHQ-9, for example, target personas scored 6.9 at mild, 12.8 at moderate, and 21.0 at severe levels---crossing the clinical cutoff ($\geq$10) between mild and moderate severity. Similarly, PTSD personas scored 7.3, 33.4, and 66.0 on the PCL-5, exceeding the clinical cutoff ($\geq$33) at moderate severity and nearly doubling it at severe.

The triviality of generating plausible counterfeits illustrates the severity of this threat to mental health research. Given that a standard commercial LLM can semi-autonomously produce clinically differentiated profiles with zero specialized configuration, the economic incentives for survey fraud are now sharply asymmetric: fabricating a large, threshold-crossing clinical cohort costs almost nothing relative to the high costs of genuine clinical data collection. Our findings do not demonstrate that such contamination is currently occurring at scale, nor do we claim that our simulation exercise captures the full complexity of authentic clinical presentations. Rather, we highlight that the technical barrier to generating plausible clinical data is remarkably low, and that the field's current defenses may not be designed to address this class of threat. The patterns of diagnostic specificity and severity sensitivity shown here were achieved with prompts that contained no information about instrument content, scoring rubrics, or clinical cutoffs. Researchers should consider implementing identity verification beyond attention checks, triangulating self-report with behavioral or physiological measures, avoiding diagnosis-specific language in recruitment materials, and treating unverified online panel data with caution\cite{quagan_navigating_2021}. Developing standardized detection methods calibrated to clinical response patterns is therefore an urgent priority.

% ============================================================
% Table 1: Diagnostic Specificity and Severity Monotonicity
% ============================================================

\begin{table}[ht]
\caption{Diagnostic Specificity and Severity Monotonicity Across Clinical Instruments}
\label{tab:main-effects}
\centering
\small
\begin{tabular}{lccccccc}
\toprule
 & \multicolumn{3}{c}{Diagnostic Specificity\textsuperscript{a}} & & \multicolumn{3}{c}{Severity Monotonicity\textsuperscript{b}} \\
\cmidrule(lr){2-4} \cmidrule(lr){6-8}
Instrument & $b$ & 95\% CI & $P$ Value & & $b$ & 95\% CI & $P$ Value \\
\midrule
PHQ-9  & 0.13  & [$-$2.25, 2.51]  & .915 & & 7.06  & [6.18, 7.94]  & $<$.001 \\
GAD-7  & $-$0.76 & [$-$2.21, 0.69] & .306 & & 6.91  & [6.49, 7.33]  & $<$.001 \\
OCI-R  & 15.76 & [12.20, 19.31]  & $<$.001 & & 18.32 & [16.90, 19.75] & $<$.001 \\
PCL-5  & 14.63 & [10.09, 19.17]  & $<$.001 & & 29.33 & [28.29, 30.37] & $<$.001 \\
MDQ    & 9.24  & [8.42, 10.07]   & $<$.001 & & 3.56  & [2.75, 4.38]  & $<$.001 \\
PQ-16  & 6.15  & [3.02, 9.28]    & $<$.001 & & 4.84  & [3.97, 5.70]  & $<$.001 \\
EDE-Q  & 2.95  & [2.26, 3.64]    & $<$.001 & & 1.70  & [1.26, 2.15]  & $<$.001 \\
\bottomrule
\end{tabular}
\begin{flushleft}
\footnotesize
Abbreviations: EDE-Q, Eating Disorder Examination Questionnaire; GAD-7, Generalized Anxiety Disorder--7; MDQ, Mood Disorder Questionnaire; OCI-R, Obsessive-Compulsive Inventory--Revised; PCL-5, PTSD Checklist for DSM-5; PHQ-9, Patient Health Questionnaire--9; PQ-16, Prodromal Questionnaire--16.
\textsuperscript{a}Regression coefficient ($b$) for the target indicator (diagnosis-congruent = 1 vs nontarget = 0) predicting instrument score, with cluster-robust standard errors by diagnosis.
\textsuperscript{b}Regression coefficient ($b$) for ordinal severity (1 = mild, 2 = moderate, 3 = severe) predicting instrument score within target groups, with cluster-robust standard errors by diagnosis. $N = 2{,}106$ synthetic personas. The dependent variable is the total instrument score (EDE-Q uses the global mean score, range 0--6). All $P$ values are two-sided.
\end{flushleft}
\end{table}

% ============================================================
% Figure
% ============================================================

\begin{figure}[p]
  \centering
  \includegraphics[width=\textwidth,height=0.85\textheight,keepaspectratio]{figures/figure1_new.pdf}
  \caption{Synthetic respondent scores across seven validated clinical instruments. Each panel shows the mean score (95\% CI) for clinically relevant diagnostic groups and an averaged clinical control group (all nontarget diagnoses) at three severity levels (light: mild; medium: moderate; dark: severe). Dashed lines indicate published clinical cutoff scores: PHQ-9 $\geq$ 10; GAD-7 $\geq$ 10; OCI-R $\geq$ 21; PCL-5 $\geq$ 33; MDQ $\geq$ 7; PQ-16 $\geq$ 6; EDE-Q $\geq$ 2.3. EDE-Q scores reflect the global mean (0--6 scale); all other instruments use summed totals. $N = 2{,}106$ synthetic personas generated by Google Gemini 2.0 Flash.}
  \label{fig:1}
\end{figure}

% ============================================================
% Online Methods (not counted in 2,000 word limit)
% ============================================================
\newpage
\section*{Online Methods}

\subsection*{Synthetic Persona Generation}

We generated 2,106 unique synthetic personas by taking the Cartesian product of 13 clinical diagnostic profiles, three severity levels (mild, moderate, severe), three age groups (25, 50, 75 years), two genders (male, female), three education levels (high school, bachelor's degree, master's degree), and three political affiliations (Democrat, Republican, Independent). Each clinical profile specified a DSM-informed diagnosis and characteristic symptoms. The 13 diagnoses spanned mood disorders (major depressive disorder, bipolar I disorder, alcohol-induced mood disorder), anxiety disorders (generalized anxiety disorder, posttraumatic stress disorder), obsessive-compulsive disorders, psychotic disorders (schizophrenia), eating disorders (anorexia nervosa, bulimia nervosa, binge eating disorder), and neurodegenerative conditions (Alzheimer disease, Parkinson disease dementia, frontotemporal dementia).

\subsection*{Clinical Assessment Batteries}

Each persona completed seven clinical instruments:
(1) the Patient Health Questionnaire--9 (PHQ-9; nine items, 0--3 scale)\cite{kroenke2001phq9};
(2) the Generalized Anxiety Disorder--7 (GAD-7; seven items, 0--3 scale)\cite{spitzer2006gad7};
(3) the Obsessive-Compulsive Inventory--Revised (OCI-R; 18 items, 0--4 scale)\cite{foa2002ocir};
(4) the PTSD Checklist for DSM-5 (PCL-5; 20 items, 0--4 scale)\cite{blevins015};
(5) the Mood Disorder Questionnaire (MDQ; 13 yes/no items)\cite{hirschfeld2000mdq};
(6) the Prodromal Questionnaire--16 (PQ-16; 16 true/false items)\cite{ising2012pq16};
and (7) the Eating Disorder Examination Questionnaire (EDE-Q 6.0; 28 items, mixed frequency and attitudinal scales)\cite{fairburn2008edeq}.
All items and response options were implemented as specified in the published instruments.

\subsection*{Simulation Procedure}

Simulations were conducted using the Google Gemini 2.0 Flash API. Procedures were adapted from \cite{westwood2025}. For each persona, the model received a system prompt containing: (1) the persona profile (diagnosis, symptoms, severity, demographics); (2) prior responses from previous batteries, ensuring cross-instrument consistency; (3) the current battery's questions with response options; and (4) instructions to respond in character, maintain consistency, and calibrate language to the persona's education level. No adversarial prompt engineering or fine-tuning was employed. The prompt contained no information about instrument content, scoring rubrics, or clinical cutoffs.

\subsection*{Scoring and Statistical Analysis}

Responses were mapped to scores using standard scoring rubrics for each instrument. We assessed diagnostic specificity using a linear regression of scores on a target indicator (diagnosis-congruent vs nontarget) per battery, with cluster-robust standard errors. Severity monotonicity was tested by regressing scores on ordinal severity within target groups.

% ============================================================
% Acknowledgments
% ============================================================
\section*{Acknowledgments}

\subsection*{Conflict of Interest Disclosures}
Dr Berner is a scientific advisor to Juniver, Ltd. The remaining authors report no conflicts of interest.

\subsection*{Data and Code Availability}

All simulation code, scoring scripts, and raw data are available at \url{https://github.com/kiante-fernandez/LLM-clinical}.

% ============================================================
% References
% ============================================================
\bibliography{references}

\end{document}
